% GENERATED FILE. Edit canonical lessons, then run npm run generate:books.
% canonical-source-hash: fnv1a64:d6947030e9034782
% canonical-lessons: FA-C457-raste, FA-C457-hojre, FA-C457-bazarche, FA-C457-timche, FA-C457-karvansara

\chapter{In the bazaar}
\label{ch:fa-in-the-bazaar}
\hlchaptermodality{457}

\begin{chapteropening}
\textbf{By the end of this chapter:} \emph{I can find my way round the bazaar: a row of shops, a bazaar shop, a small market, a covered hall and a caravanserai.}

Complete the last lesson of chapter 457: I can find my way round the bazaar: a row of shops, a bazaar shop, a small market, a covered hall and a caravanserai.
\end{chapteropening}

\section[\texorpdfstring{\fa{راسته} (râste)}{راسته (râste)}]{\fa{راسته} (râste) — a row (of the bazaar given to one trade)}

\label{lesson:FA-C457-raste}



\begin{quote}
Before the new one: say the Persian for an estate agency, then the Persian for to get in touch.
\end{quote}

\subsection*{You'll want to know: \fa{راسته}}
\textbf{\fa{راسته}} (\emph{râste}) — \textquotedblleft{}a row (of the bazaar given to one trade)\textquotedblright{}.

\begin{quote}
\fa{در} \fa{راسته} \fa{فرش} \fa{فروشها} \fa{قدم} \fa{زدیم}.
\end{quote}

\begin{quote}
\emph{Dar râste-ye farsh-forush-hâ qadam zadim.}
\end{quote}

\begin{quote}
We walked along the carpet sellers' row.
\end{quote}

Each \emph{râste} of an old bazaar belongs to one trade: gold, copper, cloth, spices.

\begin{cousinweb}
\emph{râste} is a native word, which EtymDB lists; it is built on \emph{râst} 'straight', usually traced to Middle Persian \emph{rāst}: a straight line of shops.
\end{cousinweb}

\begin{grammarlens}[title={What you've built}]
One more word for telling, explaining and describing at B1.
\end{grammarlens}

\subsection*{Your turn}
\begin{itemize}
\raggedright
  \item \emph{Say it:} \emph{râste}
  \item \emph{Say it:} \emph{râste}, once more
  \item \emph{Say it:} a sentence of your own with \emph{râste}
  \item \emph{Recall:} say the Persian for an estate agency, then the Persian for to get in touch, then say \emph{râste} again
\end{itemize}

\begin{tcolorbox}[breakable,colback=teal!4,colframe=teal!35!black,title={Before you move on}]
What does \fa{راسته} mean? (\textquotedblleft{}A row (of the bazaar given to one trade)\textquotedblright{}.) Say it once more.
\end{tcolorbox}
\section[\texorpdfstring{\fa{حجره} (hojre)}{حجره (hojre)}]{\fa{حجره} (hojre) — a shop or office in the bazaar; a small room}

\label{lesson:FA-C457-hojre}



\begin{quote}
Before the new one: say the Persian for to get in touch, then the Persian for a row.
\end{quote}

\subsection*{You'll want to know: \fa{حجره}}
\textbf{\fa{حجره}} (\emph{hojre}) — \textquotedblleft{}a shop or office in the bazaar; a small room\textquotedblright{}.

\begin{quote}
\fa{پدربزرگم} \fa{در} \fa{بازار} \fa{یک} \fa{حجره} \fa{داشت}.
\end{quote}

\begin{quote}
\emph{Pedarbozorgam dar bâzâr yek hojre dâsht.}
\end{quote}

\begin{quote}
My grandfather had a shop in the bazaar.
\end{quote}

Arabic plural \emph{hojarât}. Also a student's cell in an old school.

\begin{cousinweb}
\emph{hojre} was borrowed from Arabic \emph{ḥujra}, which EtymDB lists; it is usually glossed 'a room, a chamber'.
\end{cousinweb}

\begin{grammarlens}[title={What you've built}]
One more word for telling, explaining and describing at B1.
\end{grammarlens}

\subsection*{Your turn}
\begin{itemize}
\raggedright
  \item \emph{Say it:} \emph{hojre}
  \item \emph{Say it:} \emph{hojre}, once more
  \item \emph{Say it:} a sentence of your own with \emph{hojre}
  \item \emph{Recall:} say the Persian for to get in touch, then the Persian for a row, then say \emph{hojre} again
\end{itemize}

\begin{tcolorbox}[breakable,colback=teal!4,colframe=teal!35!black,title={Before you move on}]
What does \fa{حجره} mean? (\textquotedblleft{}A shop or office in the bazaar; a small room\textquotedblright{}.) Say it once more.
\end{tcolorbox}
\section[\texorpdfstring{\fa{بازارچه} (bâzârche)}{بازارچه (bâzârche)}]{\fa{بازارچه} (bâzârche) — a small bazaar, a local market}

\label{lesson:FA-C457-bazarche}



\begin{quote}
Before the new one: say the Persian for a row, then the Persian for a shop or office in the bazaar.
\end{quote}

\subsection*{You'll want to know: \fa{بازارچه}}
\textbf{\fa{بازارچه}} (\emph{bâzârche}) — \textquotedblleft{}a small bazaar, a local market\textquotedblright{}.

\begin{quote}
\fa{نزدیک} \fa{خانه} \fa{ما} \fa{یک} \fa{بازارچه} \fa{قدیمی} \fa{هست}.
\end{quote}

\begin{quote}
\emph{Nazdik-e khâne-ye mâ yek bâzârche-ye qadimi hast.}
\end{quote}

\begin{quote}
There is an old little market near our house.
\end{quote}

\emph{bâzâr} + \emph{-che} 'little'.

\begin{cousinweb}
\emph{bâzârche} is built from \emph{bâzâr} 'market', inherited from Middle Persian \emph{wāzār} (EtymDB), and the native ending \emph{-che} 'little'.
\end{cousinweb}

\begin{grammarlens}[title={What you've built}]
One more word for telling, explaining and describing at B1.
\end{grammarlens}

\subsection*{Your turn}
\begin{itemize}
\raggedright
  \item \emph{Say it:} \emph{bâzârche}
  \item \emph{Say it:} \emph{bâzârche}, once more
  \item \emph{Say it:} a sentence of your own with \emph{bâzârche}
  \item \emph{Recall:} say the Persian for a row, then the Persian for a shop or office in the bazaar, then say \emph{bâzârche} again
\end{itemize}

\begin{tcolorbox}[breakable,colback=teal!4,colframe=teal!35!black,title={Before you move on}]
What does \fa{بازارچه} mean? (\textquotedblleft{}A small bazaar, a local market\textquotedblright{}.) Say it once more.
\end{tcolorbox}
\section[\texorpdfstring{\fa{تیمچه} (timche)}{تیمچه (timche)}]{\fa{تیمچه} (timche) — a covered hall in the bazaar, an arcade}

\label{lesson:FA-C457-timche}



\begin{quote}
Before the new one: say the Persian for a shop or office in the bazaar, then the Persian for a small bazaar.
\end{quote}

\subsection*{You'll want to know: \fa{تیمچه}}
\textbf{\fa{تیمچه}} (\emph{timche}) — \textquotedblleft{}a covered hall in the bazaar, an arcade\textquotedblright{}.

\begin{quote}
\fa{فرشها} \fa{را} \fa{در} \fa{تیمچه} \fa{میفروشند}.
\end{quote}

\begin{quote}
\emph{Farsh-hâ râ dar timche miforushand.}
\end{quote}

\begin{quote}
They sell the carpets in the covered hall.
\end{quote}

A \emph{timche} is often domed and finely built, a hall for wholesale trade.

\begin{cousinweb}
\emph{timche} is built from \emph{tim} 'a trading hall, a caravanserai' and the ending \emph{-che} 'little'. The origin of \emph{tim} is not settled, and EtymDB lists neither word.
\end{cousinweb}

\begin{grammarlens}[title={What you've built}]
One more word for telling, explaining and describing at B1.
\end{grammarlens}

\subsection*{Your turn}
\begin{itemize}
\raggedright
  \item \emph{Say it:} \emph{timche}
  \item \emph{Say it:} \emph{timche}, once more
  \item \emph{Say it:} a sentence of your own with \emph{timche}
  \item \emph{Recall:} say the Persian for a shop or office in the bazaar, then the Persian for a small bazaar, then say \emph{timche} again
\end{itemize}

\begin{tcolorbox}[breakable,colback=teal!4,colframe=teal!35!black,title={Before you move on}]
What does \fa{تیمچه} mean? (\textquotedblleft{}A covered hall in the bazaar, an arcade\textquotedblright{}.) Say it once more.
\end{tcolorbox}
\section[\texorpdfstring{\fa{کاروانسرا} (kârvânsarâ)}{کاروانسرا (kârvânsarâ)}]{\fa{کاروانسرا} (kârvânsarâ) — a caravanserai, a roadside inn for caravans}

\label{lesson:FA-C457-karvansara}



\begin{quote}
Before the new one: say the Persian for a small bazaar, then the Persian for a covered hall in the bazaar.
\end{quote}

\subsection*{You'll want to know: \fa{کاروانسرا}}
\textbf{\fa{کاروانسرا}} (\emph{kârvânsarâ}) — \textquotedblleft{}a caravanserai, a roadside inn for caravans\textquotedblright{}.

\begin{quote}
\fa{مسافران} \fa{قدیم} \fa{شب} \fa{را} \fa{در} \fa{کاروانسرا} \fa{میگذراندند}.
\end{quote}

\begin{quote}
\emph{Mosâferân-e qadim shab râ dar kârvânsarâ migozarândand.}
\end{quote}

\begin{quote}
Travellers in the old days spent the night in a caravanserai.
\end{quote}

\emph{kârvân} 'caravan' + \emph{sarâ} 'house'. English \emph{caravanserai} comes from the Persian word.

\begin{cousinweb}
\emph{kârvânsarâ} is a native compound of \emph{kârvân} 'caravan', which EtymDB traces to Middle Persian \emph{kārawān} 'caravan; an army on the march', and \emph{sarâ} 'house, palace'. EtymDB gives the Persian word as the source of English \emph{caravanserai}.
\end{cousinweb}

\begin{grammarlens}[title={What you've built}]
One more word for telling, explaining and describing at B1.
\end{grammarlens}

\subsection*{Your turn}
\begin{itemize}
\raggedright
  \item \emph{Say it:} \emph{kârvânsarâ}
  \item \emph{Say it:} \emph{kârvânsarâ}, once more
  \item \emph{Say it:} a sentence of your own with \emph{kârvânsarâ}
  \item \emph{Recall:} say the Persian for a small bazaar, then the Persian for a covered hall in the bazaar, then say \emph{kârvânsarâ} again
\end{itemize}

\begin{tcolorbox}[breakable,colback=teal!4,colframe=teal!35!black,title={Before you move on}]
What does \fa{کاروانسرا} mean? (\textquotedblleft{}A caravanserai, a roadside inn for caravans\textquotedblright{}.) Say it once more.
\end{tcolorbox}
